\chapter{Free to Choose}

We do not begin this lecture with a fresh definition of libertarianism. We begin where Michael Sandel begins: with unfinished business from John Stuart Mill. These companion notes follow that order closely, because the turn to libertarianism is not a change of topic but the next move once Mill's utilitarian repair begins to wobble. The lecture first reopens the problem of higher pleasures, then presses Mill's defense of rights, then turns to a stronger theory of rights, and only after that asks what such a theory would require of the state.

\begin{remark}
No validated mathematical frame survives for this lecture. The displayed formulas below are therefore transcript-backed shorthand or cautious editorial reconstructions of structures that Sandel states in prose.
\end{remark}

\section{From Mill's Repair to the Need for Stronger Rights}

Sandel opens with a recap rather than a reset. Mill had tried to answer Bentham's critics by showing that utilitarianism can distinguish higher from lower pleasures. The classroom experiment with \emph{The Simpsons} and Shakespeare, however, produced a result that exposes the difficulty at once: many students reported preferring the one while still judging the other to be higher or worthier. The pressure point can be written compactly as
\begin{equation}
\text{Preferred} \neq \text{Worthier}.
\end{equation}
This is not Mill's own formula. It is simply a clean way to mark the dilemma Sandel wants us to feel before the lecture moves on.

From there the lecture turns to Mill's defense of justice in Chapter V of \emph{Utilitarianism}. Mill wants to say that rights are worthy of special respect, indeed that justice is the most sacred and most binding part of morality. Sandel grants, for the sake of argument, the stable core of Mill's long-run reply:
\begin{equation}
\begin{aligned}
\text{Respect rights}
&\Longrightarrow \text{better social consequences} \\
&\quad \text{in the long run}.
\end{aligned}
\end{equation}
The lecture immediately asks us not to settle too quickly for that answer. Sandel separates two objections that are easy to blur together.

The first objection is about hard cases. If violating a right in some exceptional case really would make people better off even in the long run, would utilitarianism then permit the violation? If so, rights remain contingent on outcome.

The second objection is deeper. Even if Mill's long-run calculus always worked as advertised, would that be the right reason to respect persons? The organ-transplant case is introduced precisely to force this question. The doctor may refrain from taking the healthy patient's organs because public trust would collapse and the long-run consequences would be bad. But Sandel asks whether there is not also, and more fundamentally, a reason of a different kind.

We can display the contrast Sandel is drawing in the organ case as follows:
\begin{align}
\text{Do not use the patient}
&\Longrightarrow \text{avoid bad} \\
&\quad \text{long-run consequences}, \\
\text{Do not use the patient}
&\Longrightarrow \text{respect the person} \\
&\quad \text{as an individual}.
\end{align}
The lecture's question is whether the first line can ever be enough by itself.

\subsection{Question \& Answer}

Even if respecting rights maximizes welfare in the long run, is that the right reason to respect persons?

Sandel's answer is that it is not enough. A rule may work well and still fail to name the proper moral reason. The point of the organ-transplant case is not only that utilitarianism may break down in difficult cases; it is that a merely instrumental defense of rights seems to miss what is morally weighty about rights in the first place. Once that gap opens, the lecture needs a stronger account of rights than Mill can comfortably provide from within utilitarianism.

\section{Libertarianism as a Strong Theory of Rights}

The transition now becomes explicit. If Mill's defense of rights seems to lean on ideas such as dignity or respect that are not themselves utilitarian, then we should ask whether there are stronger theories of rights that put those ideas first. Sandel says, in effect: today we turn to one of them.

The defining claim is that individuals matter not merely as containers of utility, not merely as locations where satisfactions can be stored and summed, but as separate beings with separate lives. That is why strong theories of rights resist the utilitarian habit of thinking about law and justice by aggregation alone.

Libertarianism enters at exactly this point. It takes liberty to be the fundamental right. Because we are separate persons, we are not available for whatever use society might find advantageous. In schematic form, the lecture pushes toward the following claim:
\begin{equation}
\begin{aligned}
\text{SelfOwnership}(p)
&\Longrightarrow \text{strong side-constraints} \\
&\quad \text{on using }p.
\end{aligned}
\end{equation}
Again, this notation is editorial. Sandel states the idea verbally. But the implication is faithful to the role the idea plays in the lecture.

This is also the setting for Nozick's formulation that rights are so strong and far-reaching that they raise the question of what, if anything, the state may do. The lecture does not leave the point at the level of slogan. As soon as liberty is introduced, Sandel asks what follows institutionally.

\section{The Minimal State and Its Three Prohibitions}

Sandel's next move is crisp and numbered. If liberty is the basic right, then the state is barred from three kinds of action that modern governments commonly take to be legitimate.

\begin{enumerate}
\item Paternalist legislation: laws that protect people from themselves.
\item Morals legislation: laws that enforce virtue or give expression to a collective moral code.
\item Redistribution: taxation or policy that transfers income or wealth from the rich to the poor.
\end{enumerate}

The lecture develops these in order. Seatbelt and motorcycle-helmet laws are the first examples, because they make the paternalist form of coercion easy to see. Laws against consensual sexual intimacy are then introduced as the paradigm of morals legislation. Redistribution comes third and receives the most attention, because it exposes the sharpest dispute about coercion, property, and the use of persons.

A compact editorial summary is
\begin{equation}
S_{\mathrm{forbidden}}=\{\text{paternalism},\ \text{morals legislation},\ \text{redistribution}\}.
\end{equation}
By contrast, the minimal state is allowed only a narrow range of functions:
\begin{equation}
S_{\min}=
\left\{
\begin{aligned}
&\text{defense},\ \text{police},\ \text{courts}, \\
&\text{contract/property enforcement}
\end{aligned}
\right\}.
\end{equation}
These set-valued expressions are note-writer's shorthand, not lecture notation. What matters is the shape of the doctrine as Sandel unfolds it: the state may protect against force and fraud, but it may not coerce us for our own good, for the sake of virtue, or for distributive equality.

Having drawn those boundaries, the lecture concentrates on the third prohibition. Sandel wants not merely to announce that redistribution is wrong on libertarian grounds, but to show what has to be true for that conclusion to follow.

\section{Redistribution, Historical Justice, and the Entitlement View}

At this stage Sandel makes the issue concrete by asking about the distribution of wealth in the United States. The country is highly unequal. Is that fact enough to tell us whether the distribution is just or unjust? The libertarian answer is no. We cannot read justice directly off the pattern.

Let \(D\) denote a distribution of holdings and \(J(D)\) its justice-status. Then the lecture's historical test can be summarized as
\begin{equation}
\begin{aligned}
J(D)\ &\text{depends on the history of acquisition} \\
&\text{and transfer, not on pattern alone}.
\end{aligned}
\end{equation}
That sentence is almost the heart of the entitlement view. To judge a distribution we need to know how it arose.

Sandel states two principles. First comes justice in acquisition: were the original holdings obtained fairly, or were land, factories, goods, or other productive assets stolen or seized? Second comes justice in transfer: did later holdings arise through free consent, through buying, selling, trading, and exchange on the market?

A compact reconstruction is
\begin{equation}
J(D)=J_{\mathrm{acq}} \land J_{\mathrm{tr}},
\end{equation}
where \(J_{\mathrm{acq}}\) records justice in acquisition and \(J_{\mathrm{tr}}\) justice in transfer. The formula is editorial, but it captures the lecture's logic without adding machinery the lecture does not use.

\paragraph{How the test works.}
The entitlement view proceeds step by step.
\begin{enumerate}
\item Begin with an observed distribution \(D\).
\item Refuse to judge \(D\) from its pattern alone.
\item Ask whether the relevant initial holdings were acquired justly.
\item Ask whether later holdings arose through voluntary transfer.
\item If both conditions are met, treat the resulting distribution as just.
\item If either condition fails, the resulting distribution is unjust.
\end{enumerate}

\subsection{Question \& Answer}

Can we judge a distribution by the end-state alone, or only by the history of acquisition and transfer?

The lecture's answer is that pattern is not enough. Libertarian justice is historical. A highly unequal distribution may be just if it arose from just acquisition and free exchange. A more equal distribution may still be unjust if it was produced by coercion, confiscation, or fraud. Sandel needs this historical turn before he can make redistribution the live issue, because otherwise the argument would be swallowed immediately by egalitarian intuition.

\section{Bill Gates, Michael Jordan, and the First Round of Objections}

Sandel now says, in effect, let us make it concrete. He turns from the abstract structure of entitlement to two examples that give the lecture its numerical backbone: Bill Gates and Michael Jordan.

The Gates example is built to activate the utilitarian impulse. Gates is presented as so rich that ordinary scales stop working. Sandel uses the lecture-era Lincoln Bedroom joke to dramatize magnitude: Gates could afford to stay there every night for the next \(66{,}000\) years. He then turns to effective earning rate. If we spread Gates's accumulated wealth over long working days since the founding of Microsoft, the lecture says his rate of pay comes out to more than \(\$150\) per second. Hence the joke that if he saw a \(\$100\) bill on the street, it would not be worth stopping to pick it up.

That is exactly the setup Sandel wants. Once one person is imagined as scarcely noticing the loss, utilitarianism pushes strongly toward redistribution for urgent needs at the bottom. The argument is not subtle:
\begin{enumerate}
\item A very small loss to Gates produces negligible pain.
\item The same amount, distributed to those in serious need, produces large gains in welfare.
\item Therefore a utilitarian would redistribute at once.
\end{enumerate}
Sandel immediately checks that conclusion by returning to the libertarian premise: if the holdings are justly his, then the gain in aggregate welfare does not by itself settle the moral question.

The Jordan example makes the issue even sharper because the arithmetic is explicit:
\begin{equation}
31\,\text{million}+47\,\text{million}=78\,\text{million}.
\end{equation}
That is, \(\$31\) million in salary plus \(\$47\) million in endorsements. Sandel then asks us to imagine taking roughly one third of that total for food, health care, housing, and education for the poor. The implied arithmetic is
\begin{equation}
\frac{1}{3}\cdot 78\,\text{million}=26\,\text{million}.
\end{equation}

\paragraph{Worked example.}
The lecture's contrast can be written in four short steps.
\begin{enumerate}
\item Begin with Jordan's annual income: \(78\) million dollars.
\item Take one third: \(26\) million dollars.
\item A utilitarian says: Jordan scarcely notices the loss, while the poor gain substantially.
\item A libertarian says: if Jordan is entitled to the income, the gain in welfare does not cancel the wrong of coercive taking.
\end{enumerate}

Now the classroom pressure begins. The first objection is that wealth like Jordan's is not simply the product of isolated effort. The successful, one student argues, have received a larger gift from society and therefore owe more back. A second objection worries that absent some redistribution, there can be no genuine equality of opportunity for those who begin far behind.

The libertarian side answers first with analogy. Joe's skateboard example asks us to imagine one person with a collection everyone else suddenly wants. If the others enter his house and take them, that is unjust even if the result is more equal. Sandel presses Joe: does a \(33\%\) tax on Jordan for good causes amount to theft? Joe answers yes, though he allows that in an extreme case an injustice may have to be condoned.

The reply to Joe then changes the register again. It is said that this is not like taking \(99\) of \(100\) skateboards; it is more like taking a portion of a surplus so large that its owner cannot meaningfully use it all. And if the state never redistributes, some citizens may be trapped so low that equal opportunity becomes hollow. The lecture leaves this first round unresolved, because Sandel wants to show that Nozick's sharper argument does not stop at theft.

\section{From Taxation to Forced Labor to Self-Possession}

At this point Sandel makes the decisive escalation. Joe had called redistribution theft. Nozick, we are told, goes one step further. Taxation does not merely take an external object. It takes earnings, and earnings are the fruits of labor. The argument therefore unfolds in a chain:
\begin{align}
\text{taxation of earnings}
&\to \text{taking fruits of labor} \\
&\to \text{claim on labor} \\
&\to \text{forced labor} \\
&\to \text{slavery}.
\end{align}

The lecture insists on the intermediate steps. If the state may claim the product of my labor, then it has, in effect, asserted a title to some portion of the labor that produced that product. If it has a title to some portion of my labor, then it is laying claim to some portion of my life, time, and effort. That is why Nozick says taxation for redistribution is morally equivalent to forced labor.

We can state the structure more explicitly as a derivation.
\begin{enumerate}
\item Taxation takes part of a person's earnings.
\item Earnings are the fruits of labor.
\item Taking the fruits of labor is morally akin to claiming part of the labor itself.
\item Claiming part of a person's labor is forced labor.
\item Forced labor means that another party is, in part, owner of the person.
\item Partial ownership of the person is slavery, or a denial of self-possession.
\end{enumerate}

That is the point at which Sandel says the stakes are very high. Once the argument is stated in this form, redistribution can no longer be treated as a merely technical dispute about fiscal policy. The dispute has become a dispute about whether the political community may own any share in the person.

The hidden premise is then brought to the surface:
\begin{equation}
\begin{aligned}
\text{SelfOwnership}(p)
&\Longrightarrow \text{the state has no title} \\
&\quad \text{to the labor of }p.
\end{aligned}
\end{equation}
This is also why Sandel deliberately links the argument back to the organ-transplant case. If we own ourselves, then the patient in the transplant case is not available for social use; and for the same reason, persons are not available for paternalist regulation, for morals legislation, or for redistributive taking in the name of the common good. One principle is meant to govern all of these cases.

\subsection{Question \& Answer}

Is taxation for redistribution really morally equivalent to forced labor?

Sandel does not ask the question in order to silence objection. He asks it in order to locate the exact point of dispute. If we reject the libertarian conclusion, we must break into the chain somewhere: perhaps by denying that taking earnings is equivalent to taking labor, perhaps by denying that claims on labor amount to ownership, or perhaps by questioning the deeper premise of self-possession itself. The lecture's achievement here is to force the disagreement into clear form.

\section{Democracy, the Minimal State Revisited, and the Turn to Locke}

The lecture now enters its functional reprise. Sandel says, in effect, that before returning to redistribution he wants to say one more word about the minimal state. This reset matters, because it shows that libertarianism is not only a theory about tax rates. It is a much broader account of what the state may and may not do.

Milton Friedman is introduced as an ally on the question of paternalism. Social Security may be prudent, even highly prudent, but on this view prudence is not enough to justify coercion. If people wish to save for retirement, they should do so; if they prefer to live lavishly now and risk a poor retirement later, that too should be their choice. The point is not that one choice is wise and the other foolish. The point is that the state may not force the choice.

The Salem Fire Corporation anecdote serves the same reprising function from a different angle. We are often told that collective goods such as police and fire protection must be publicly provided because otherwise free riders cannot be controlled. Sandel uses the private fire-company story to show how a libertarian might deny that inevitability. The company responds only to subscribers, or to fires threatening subscribers. A house can burn while the trucks stand by, because if the company extinguished every fire there would be no incentive to subscribe. Sandel does not present this as a theorem; he presents it as a stress test for our assumptions about public goods.

From there the lecture returns to the classroom in a more organized way. Team Libertarian is assembled, and Sandel states the objections in order. The poor need the money more. Taxation in a democracy is not slavery because it is authorized by consent of the governed. The successful owe a debt to society that is properly repaid by taxation.

The reply to the first objection is the stable core that survives the corrupted transcript: even if Bill Gates may choose to give voluntarily, coercive redistribution still violates property rights on the libertarian view. Julia sharpens that reply by distinguishing need from desert. Need may be urgent and morally important, but it does not by itself show entitlement to another person's holdings. That contrast can be kept in view schematically as
\begin{align}
\text{Need} &\neq \text{Desert}, \\
\text{Need help} &\not\Longrightarrow \text{entitled to another's holdings}.
\end{align}
The Katrina case is then used to press this distinction. Julia concedes the urgency of need while maintaining that need and desert are not the same category.

The second objection turns on democracy. Raul argues that property rights are enforced by democratic government, and so taxation under those rules is legitimate rather than coercive. John replies in rough numerical terms that redistribution becomes the middle deciding what the top will do for the bottom. Alex then gives the more philosophically important answer: having a tiny share in an election is not the same as exercising personal control over one's own labor and property. One may vote and still lose. Consent through a vast representative machinery is not, on this view, equivalent to self-direction.

Sandel pushes this objection hard, and the exchange improves. Must not a democrat accept majority rule? Alex's reply is not to reject democracy wholesale, but to restrict its scope. Fundamental rights should not be put to a vote. Sandel then supplies the telling analogy: if religious liberty were at stake, we would hesitate to let a majority abolish it. The libertarian therefore tries to assimilate rights of property and earnings to the same family of rights as free speech and religious liberty.

The third objection is that the successful owe a debt to society. Julia answers that the successful became successful precisely by doing something society valued highly, and that they have already been paid through the voluntary exchanges that made them rich. Sandel makes the Michael Jordan case concrete: coaches, teammates, teachers, and spectators all matter, but they have already been compensated or have already received the entertainment and satisfaction for which they paid. The repeated transcript line about society's pleasure in watching Jordan is best compressed into one clean point: on the libertarian reply, the service has already been rendered and rewarded.

At this point Victoria raises the deepest objection of the lecture. Maybe, she suggests, we do not own ourselves after all. Living in society already places limits on what we may do. One cannot simply invoke self-possession as if the presence of others changed nothing. That objection does not merely dispute one tax policy. It questions the premise that holds the entire libertarian edifice together.

\subsection{Question \& Answer}

If living in society already limits what we may do, do we really own ourselves in the strong libertarian sense?

Sandel does not answer this by returning to utilitarianism. He treats it as the decisive philosophical hinge. If self-possession fails, then the libertarian argument against redistribution loses its deepest support. But if we abandon self-possession too quickly, we risk sliding back toward a view in which persons may once again be used for aggregate welfare. The lecture therefore ends not with closure but with displacement. The argument must now be carried to a deeper source.

That source is Locke. Sandel closes by sketching the chain Nozick inherits:
\begin{equation}
\begin{aligned}
\text{self-possession}
&\to \text{ownership of labor} \\
&\to \text{mixing labor with unowned things} \\
&\to \text{private property}.
\end{aligned}
\end{equation}
The point of the chapter's ending is not to settle the matter but to show exactly why the next lecture has to be about Locke. If self-ownership is the premise, then the next task is to test that premise where it first receives its classical philosophical form.

\section{Summary}

This lecture begins with Mill and ends with Locke, but the movement is continuous. We start from Mill's attempt to repair utilitarianism, first with higher pleasures and then with rights. We then see why Sandel thinks the long-run utilitarian defense of rights is unstable: it may fail in hard cases, and even when it succeeds it may still offer the wrong kind of reason. That pressure opens the way to libertarianism as a stronger theory of rights.

Once libertarianism appears, the lecture proceeds in a deliberate sequence. Liberty is translated into a view of the state. Redistribution is then tested by a historical account of justice rather than a patterned one. Gates and Jordan make the stakes numerically vivid. Nozick pushes the argument from theft to forced labor and from forced labor to slavery. The later classroom debate then reveals where resistance gathers: need, equality of opportunity, democratic consent, debt to society, and finally the premise of self-possession itself.

So the lecture ends where it should: not with a verdict on tax policy, but with a deeper question. Do we own ourselves? Sandel leaves that question suspended and turns us toward Locke, because the answer to it will determine how much of the libertarian case can stand.
